Reinforcement learning (RL) \citep{RL_intro} has gained increasing attention in control and
transportation research, as it enables control policies to be learned through
interaction with a simulated environment \citep{EPSILON_GREEDY, busoniu2018reinforcement}.
In intelligent transportation systems, RL has been applied to freeway control
problems such as ramp metering, variable speed limits, and combined traffic
management strategies \citep{han2023leveraging}.\\
Ramp metering regulates on-ramp inflow to reduce congestion and improve
mainline throughput \citep{papageorgiou2002freeway}. Classical feedback
controllers, such as ALINEA and PI-ALINEA, use traffic measurements near the
control target location to determine the metering rate \citep{ALINEA, PI-ALINEA}.
While these methods are well established for bottlenecks close to the ramp,
distant downstream bottlenecks introduce delays between metered ramp vehicles
and their effect at the bottleneck, which can reduce controller stability. This motivates control strategies that can anticipate traffic
evolution between the on-ramp and the downstream bottleneck.\\
This work aims to reproduce the study by \cite{ORIGINAL_STUDY}, which formulates
ramp metering for a distant downstream bottleneck as a Q-learning problem with
value function approximation. In the original study, the agent observes traffic
densities at representative locations between the ramp and the bottleneck,
together with the estimated ramp demand, and selects a discrete metering rate.
An artificial neural network approximates the action-value function, allowing the
method to operate in a multidimensional continuous state space without a large
lookup table.\\
The purpose of this reproduction is to evaluate the reported performance of the
proposed RL controller and to assess the reproducibility of the experimental setup.
As highlighted by \cite{riehl2025revisiting}, transportation simulation studies
often face reproducibility challenges due to missing source code, incomplete
calibration information, and insufficiently documented modelling assumptions.
Although the original article provides the main modelling concepts and controller
structure, several implementation details required for exact reproduction are
unspecified or only available graphically. This study therefore reconstructs the
simulation and learning pipeline from the published description and documents
the assumptions required to reproduce the reported results.\\
The remainder of this paper is structured as follows. Section~\ref{sec:Summary_of_origninal_study}
summarises the original study. Section~\ref{sec:Implementation_details}
describes the reproduction methodology. Section~\ref{sec:Results_and_Discussion}
presents and discusses the reproduced results. Finally,
Section~\ref{sec:Conclusion} concludes this replication study.

%\clearpage